\documentclass[11pt]{article}
\usepackage[margin=1in]{geometry}

\usepackage{amsmath,amssymb,amsthm,mathtools}
\usepackage[T1]{fontenc}
\usepackage{hyperref}

\setlength{\parindent}{0pt}
\setlength{\parskip}{0.7\baselineskip}

\newtheorem{definition}{Definition}
\newtheorem{assumption}{Assumption}
\newtheorem{proposition}{Proposition}
\newtheorem{theorem}{Theorem}
\newtheorem{corollary}{Corollary}
\newtheorem{remark}{Remark}

\newcommand{\Warm}{\textnormal{\textsc{WarmStartExactJh}}}
\newcommand{\Act}{\mathcal{A}}
\newcommand{\Ospace}{\mathcal{O}}
\newcommand{\Rspace}{\mathcal{R}}
\newcommand{\Hist}{\mathcal{H}}
\newcommand{\Teach}{\mathcal{D}_{\mathrm{teach}}}
\newcommand{\TV}{\mathrm{TV}}
\newcommand{\Jh}{\mathcal{J}_H}
\newcommand{\Bridge}{\mathfrak{B}}
\newcommand{\Lab}{\mathcal{L}_H}
\newcommand{\Vlab}{V_{\Lab}}
\newcommand{\PT}{P_T}
\newcommand{\QT}{Q_T}
\newcommand{\Qxi}{Q_{\xi,H}}

\title{\Warm:\\Standalone Exact Finite-Horizon Warm-Start Agent Specification}
\author{Noah Cashin}
\date{\today}

\begin{document}
\maketitle

\section{Status and Scope}
\label{sec:status}

This document is the normative standalone specification for the
\Warm{} controller family implemented by
\[
\texttt{controller.kind = "aiqi\_warmstart\_exact\_jh"}.
\]
It specifies the mathematical model, canonical teacher artifact, offline
teacher export, online trace refresh, and implementation contract. It is
intentionally independent of \texttt{docs/infotheory-tuner-v1.tex}; the tuner
use case is an instantiation of the model, not the definition of the model.

The core object is a finite-action, finite-observation, finite-label
controller over a certified same-task bridge. The bridge defines admissible
actions, observations, one-step rewards, finite-horizon labels, label decoding,
and task identity. \Warm{} learns a conditional return-label law
\[
\widehat p(\ell\mid h,a),
\]
not an action-cloning policy \(\pi_T(a\mid h)\). This is the structural sense
in which it is AIQI-like: it performs direct induction over action-conditional
return distributions rather than imitating the teacher's realized action
sequence.

A teacher is any controller or policy that can produce realized transitions
under the same compiled bridge contract. Its family name is not part of
admissibility: MC-AIXI, discounted AIQI, \Warm{} itself, a hand-coded policy,
a random policy, or a future Infotheory controller are all valid teachers if
their traces satisfy the contract below.

The Rust/controller name \Warm{} is historical for the tuner exact-\(\Jh\)
instantiation. Mathematically, the core is bridge-parametric exact return-law
warmstart. In the direct integer-reward bridge, labels decode to \(H\)-step
reward sums. In the tuner bridge, labels decode to exact \(\Jh\) incumbent
improvement. This document preserves the tuner theorem boundary by treating
that bridge as a separately certified instantiation.

\section{Compiled Task Contract}
\label{sec:task}

\begin{definition}[Certified same-task bridge]
A certified same-task bridge is a tuple
\[
\Bridge=(\Act,\Ospace,\Rspace,\Lab,\Vlab,\Omega_R,\Omega_{\Lab},
\operatorname{Label}_H,\sigma)
\]
where \(\Act\), \(\Ospace\), and \(\Rspace\) are finite action,
observation, and one-step reward domains; \(\Lab\) is a finite label
alphabet; \(\Vlab:\Lab\to\mathbb{R}\) is the exact semantic value decoder;
\(\Omega_R\) and \(\Omega_{\Lab}\) are injective finite encoders for rewards
and labels; \(\operatorname{Label}_H\) deterministically maps every realized
same-task \(H\)-window to a label in \(\Lab\); and \(\sigma\) is the bridge
task fingerprint.
\end{definition}

\Warm{} only requires the bridge to provide finite, same-task, exactly
decodable labels. The direct standalone bridge uses
\[
\operatorname{Label}_H((a_j,o_j,r_j)_{j=i}^{i+H-1})
=
\sum_{j=i}^{i+H-1}r_j
\]
and \(\Vlab(\ell)=\ell\). Other certified bridges may use different exact
label semantics without changing the warmstart predictor or policy.

\begin{definition}[Compiled warm-start task]
A compiled warm-start task is a tuple
\[
M=(\Bridge,H,B,N,P,\rho)
\]
with the following components:
\begin{itemize}
    \item a certified same-task bridge \(\Bridge\);
    \item a finite return horizon \(H\ge1\);
    \item an exact return-label alphabet size
    \(B=|\Lab|\ge1\);
    \item a label phase period \(N\ge H\);
    \item a finite-alphabet predictive backend \(P\) supporting strict frozen
    history conditioning;
    \item a deterministic seed \(\rho\).
\end{itemize}
\end{definition}

For the direct integer-reward standalone bridge, the return-label alphabet must
be exact for a dense nonnegative reward interval. Thus \(B-1\) must be
divisible by \(H\), and the reward domain is derived from \(H\), \(B\), and
\(b_R\) by
\[
r_{\min}=0,
\qquad
r_{\max}=\frac{B-1}{H},
\qquad
\Rspace=\{0,\dots,r_{\max}\}.
\]
The channel bit width must represent every reward in \(\Rspace\). This
choice makes the \(H\)-step return set exactly
\(\{0,\dots,B-1\}\). Slack return bins are non-conforming: they would create
labels with no exact direct-bridge return preimage and are rejected by the
compiled warm-start contract. Signed or non-identity reward channels require
an explicit future contract or an externally certified bridge such as the
tuner bridge in Section~\ref{sec:tuner-bridge}.

At step \(t\ge1\), the controller emits \(a_t\in\Act\), then receives
\((o_t,r_t)\in\Ospace\times\Rspace\). The realized history before choosing
\(a_t\) is
\[
h_{<t}=(a_1,o_1,r_1,\dots,a_{t-1},o_{t-1},r_{t-1})\in\Hist.
\]
For any start step \(i\), once the next \(H\) rewards are known, define the
exact delayed label
\[
\ell_i=\operatorname{Label}_H((a_j,o_j,r_j)_{j=i}^{i+H-1}).
\]
For the direct integer-reward bridge this is
\[
G_i^{(H)}=\sum_{j=i}^{i+H-1}r_j,
\qquad
\ell_i=G_i^{(H)}.
\]
The implementation stores direct-bridge labels as unsigned integers; the semantic value
used for action selection is always \(\Vlab(\ell_i)\).

Throughout this document, total variation uses the standard probability
metric convention
\[
\|p-q\|_{\TV}
=
\sup_A |p(A)-q(A)|
=
\frac12\sum_x |p(x)-q(x)|.
\]

The phase of a one-based step index is
\[
\phi(i)=i\bmod N\in\{0,\dots,N-1\}.
\]
Phase partitioning prevents unresolved delayed labels from being required in
the same predictor stream that is being conditioned for a current decision.

\section{Canonical Teacher Artifact}
\label{sec:teacher-artifact}

\subsection{JSON object}

The canonical teacher artifact is a versioned JSON object
\[
(\texttt{schema\_version},\texttt{contract},\texttt{traces})
\]
with schema version \(1\). Its shape is:
\begin{verbatim}
{
  "schema_version": 1,
  "contract": {
    "task_fingerprint": "...",
    "action_alphabet_size": K,
    "observation_bits": b_O,
    "observation_stream_len": L,
    "observation_key_mode": "first|last|stream_hash|full_stream",
    "observation_adapter_spec_ref": "...",
    "observation_adapter_content_crc32": "........",
    "reward_bits": b_R,
    "return_horizon": H,
    "label_phase_period": N,
    "scalar_representation": "...",
    "exact_reward_encoding_certificate": "..."
  },
  "traces": [
    {
      "transitions": [
        { "action": a, "observations": [o_0, ...], "reward": r }
      ]
    }
  ]
}
\end{verbatim}
Legacy trace arrays and field aliases such as \texttt{obs} are not part of
the schema.

The field \texttt{return\_bins} is not duplicated in the teacher artifact.
It is part of the target planner-run controller spec, and the teacher contract
is validated against that target. Thus \(B=|\Lab|\) is recovered from the
compiled target, while the artifact records the horizon and phase-period
fields needed to interpret transition windows.

For standalone planner runs, the observation adapter declaration is
\[
\texttt{direct-planner-percept-lsb-first-v1},
\]
and the scalar representation is
\[
\texttt{nonnegative-integer-i64-instantaneous-reward-v1}.
\]
The adapter CRC and reward certificate are CRC32 hex strings of the canonical
JSON declaration objects implemented in
\texttt{aixi/warmstart\_contract.rs}. They commit to the observation stream
length, observation bit width, reward bit width, LSB-first field encoding, and
identity reward encoder over the representable nonnegative reward interval.

\subsection{Admissibility}

\begin{definition}[Admissible same-task teacher artifact]
\label{def:admissible-teacher}
A teacher artifact \(\Teach\) is admissible for a compiled target task \(M\)
when:
\begin{itemize}
    \item \texttt{schema\_version} is \(1\);
    \item the contract task fingerprint equals the target fingerprint
    \(\sigma\);
    \item the contract action-alphabet size, observation bit width,
    observation stream length, observation key mode, reward bit width, return
    horizon, and label phase period equal the compiled target contract;
    \item when the target environment is not the externally validated tuner
    bridge, the standalone observation-adapter, scalar-representation, and
    exact-reward-encoding provenance fields match the canonical declarations
    for the compiled target;
    \item every transition action lies in \(\Act\), every observation cell lies
    in its declared bit range, and every reward lies in \(\Rspace\);
    \item every trace has length at least \(H\), so each trace contributes at
    least one complete delayed-label window;
    \item every complete \(H\)-window in every trace yields a label in
    \(\Lab\);
    \item the artifact contains at least one complete delayed label in total.
\end{itemize}
\end{definition}

\subsection{Task fingerprint}
\label{sec:task-fingerprint}

The standalone task fingerprint is a SHA-256 hex digest of a deterministic JSON
payload containing:
\begin{itemize}
    \item the SHA-256 of the canonical planner-run document bytes with the
    teacher dataset asset removed from both \texttt{assets} and the controller
    binding;
    \item the compiled controller kind;
    \item the compiled controller backend label;
    \item the teacher contract schema version;
    \item SHA-256 content commitments for every resolved task asset except the
    teacher dataset asset itself.
\end{itemize}
Excluding the teacher dataset asset is essential: it permits a target \Warm{}
planner-run document to name an output teacher asset that does not exist yet,
and it prevents the artifact bytes from changing the task identity they are
meant to satisfy.

\subsection{Telemetry versus teacher data}
\label{sec:telemetry}

JSONL action/percept logs are raw telemetry. They are not canonical teacher
data. The telemetry grammar consists of one JSON object per line:
\[
\{\texttt{"t"}:n,\texttt{"kind"}:\texttt{"action"},\texttt{"action"}:a,
\texttt{"provenance"}:s\}
\]
or
\[
\{\texttt{"t"}:n,\texttt{"kind"}:\texttt{"percept"},
\texttt{"observations"}:[o_0,\dots,o_{L-1}],\texttt{"reward"}:r\}.
\]
For action records, \texttt{provenance} is optional for legacy telemetry and
required for conforming warmstart experiment logs. Its allowed values are:
\[
\texttt{"greedy"},\quad
\texttt{"exploratory"}.
\]
The converter ignores provenance when constructing teacher data; provenance is
decision accounting, not trace content. If \texttt{provenance} is present, a
conforming converter must validate that it is one of the allowed labels above;
omitted provenance is accepted only for legacy telemetry.

A deterministic converter from telemetry to a teacher artifact must:
\begin{enumerate}
    \item reject malformed JSON, unknown fields needed for conversion,
    duplicate action records at the same step, and duplicate percept records
    at the same step;
    \item pair an action at step \(n\) with a percept at the same step only
    when the percept record occurs after the action record;
    \item otherwise pair the action at step \(n\) with the percept at
    step \(n+1\), supporting pre-action percept logs;
    \item reject any action for which no such post-action percept exists;
    \item validate every reconstructed transition against the target bridge
    contract;
    \item attach the target contract and reject traces shorter than \(H\);
    \item emit the canonical JSON teacher dataset above.
\end{enumerate}
This converter never reorders records to repair a broken log and never fills
missing percepts. A single JSONL trace must use one pairing convention
throughout. Mixed pre-action and post-action telemetry is non-conforming and
must be rejected rather than heuristically repaired.

Multiple teacher artifacts for the same contract are merged by deterministic
deduplicating union over exact trace content. A trace is identified for merge
purposes by its transition list itself, not by a hash, CRC, or serialized JSON
digest. Transition lists are compared structurally and lexicographically: each
transition is ordered by action, then by the observation stream, then by reward;
if one transition list is a prefix of another, the shorter trace sorts first.
Identical transition lists are deduplicated, and newly inserted traces are
placed according to this structural order. Source-controller metadata is
provenance, not part of trace identity or ordering. Hashes of trace content may
be reported as diagnostics, but they are not normative merge keys.

\section{Predictor Construction and Encoding}
\label{sec:predictor}

Let
\[
b_A=\lceil\log_2 K\rceil,\qquad b_G=\lceil\log_2 B\rceil.
\]
Action, observation, and reward integer fields are encoded LSB-first with
fixed width:
\[
\operatorname{enc}_A(a)\in\{0,1\}^{b_A},
\operatorname{enc}_O(o_j)\in\{0,1\}^{b_O},
\quad
\operatorname{enc}_R(r)\in\{0,1\}^{b_R}.
\]
Return labels use the value-monotone fixed-width code
\[
\operatorname{enc}_G(\ell)\in\{0,1\}^{b_G}
\]
given by the MSB-first natural binary representation of the ordered label
index. Thus, for power-of-two label alphabets, every depth-\(d\) label-prefix
node denotes a contiguous block of decoded return values with width at most
approximately \(2^{-d}\) of the full return range. This is intentionally
different from the LSB-first field encoding used for actions, observations,
and rewards. The separation is semantically important: changing label bit
order changes the return-predictor training stream, while changing action or
percept field order would change the bridge stream itself.

\Warm{} maintains one predictive model \(P_n\) for each phase
\(n\in\{0,\dots,N-1\}\). For a trace
\[
\tau=((a_1,o_1,r_1),\dots,(a_T,o_T,r_T)),
\]
compute labels \(\ell_i\) for \(1\le i\le T-H+1\). The phase-\(n\)
training stream for this trace is the concatenation, for \(i=1,\dots,T\), of
\[
\operatorname{enc}_A(a_i),
\quad
\begin{cases}
\operatorname{enc}_G(\ell_i), & \phi(i)=n \text{ and } i\le T-H+1,\\
\epsilon, & \text{otherwise},
\end{cases}
\quad
\operatorname{enc}_O(o_i),
\quad
\operatorname{enc}_R(r_i).
\]
The implementation commits this stream into every phase model. After each
teacher trace, the predictor's transient conditioning history is reset while
its learned state remains.

During live interaction, when the label for start step \(i\) becomes known,
only phase \(P_{\phi(i)}\) is advanced through the corresponding augmented
live stream. This keeps each phase model chronologically consistent with the
labels it is allowed to see.

\section{Action Values and Policy}
\label{sec:policy}

At the next step \(t\), the active phase is \(n=\phi(t)\). For a candidate
action \(a\), \Warm{} conditions \(P_n\) on any uncommitted live history,
pushes \(\operatorname{enc}_A(a)\) as temporary history, and queries the
bit-chain probabilities of the return-label law. The implementation uses a
memoized shared-prefix trie evaluator as the complete authoritative baseline:
one logical next-bit query is made per reached internal label-prefix node, and
the resulting prefix masses are reused for all descendant labels. In the clean
power-of-two case this is at most \(B-1\) logical label-bit probability
queries per action, instead of \(B b_G\) leaf-by-leaf bit queries. The
resulting nonnegative masses are normalized to a distribution
\[
\widehat p_t(\ell\mid h_{<t},a).
\]
If the mass sum is zero or non-finite, the implementation falls back to the
uniform distribution over labels. The estimated value is
\[
\widehat Q_H(h_{<t},a)=
\sum_{\ell\in\Lab}\Vlab(\ell)\,
\widehat p_t(\ell\mid h_{<t},a).
\]

The greedy action is
\[
\hat a_t=\min\arg\max_{a\in\Act}\widehat Q_H(h_{<t},a).
\]
An implementation may expose additional epsilon exploration during learning:
with probability \(\varepsilon_t\) it samples uniformly from \(\Act\), otherwise
it executes \(\hat a_t\). Exact greedy claims apply only when the executed
decision is greedy.

After executing \(a_t\), the controller records \((a_t,o_t,r_t)\). If this
creates a new delayed label \(\ell_{t-H+1}\), the active phase for that label
is updated before future predictions require it.

\subsection{Online learning modes}
\label{sec:online-modes}

\Warm{} has three distinct online data modes:
\begin{enumerate}
    \item within-run delayed-label update: the current agent updates a phase
    predictor as soon as a realized \(H\)-window closes;
    \item exploration-supported correction: the run deliberately executes
    non-greedy actions to improve action-conditional coverage, and exact
    greedy-action claims are disabled at those exploratory decisions;
    \item across-run refresh: a completed live trace is merged into the teacher
    artifact and a new student is rebuilt from the merged dataset.
\end{enumerate}
Within-run updates and across-run rebuilds may both be supported by an
implementation, but a single reported training event must be classified as
one or the other. The tuner executor currently reports either online delayed
label update or trace-refresh rebuild for the same self-improvement round, so
the theorem/provenance accounting is not double-counted. In the current tuner
bridge, trace refresh disables the same round's online delayed-label update
flag; standalone warmstart runtimes may still update within a run, then later
perform an explicit across-run refresh as a separate event.

Online updates improve the model only on histories and actions that are
actually covered. They do not by themselves create counterfactual labels for
untried actions. Exploration-supported correction is therefore a coverage
mechanism, not a theorem of monotone improvement.

\section{Offline Teacher Export}
\label{sec:offline-export}

The canonical CLI export form is:
\begin{verbatim}
infotheory warmstart teacher planner-run \
  --target <warmstart-planner-run> \
  --teacher <teacher-planner-run> \
  --out <teacher.json>
\end{verbatim}
The target planner run must use \texttt{aiqi\_warmstart\_exact\_jh}. The
teacher planner run may use any CLI-executable planner controller. The export
command:
\begin{enumerate}
    \item compiles both planner-run documents through the canonical spec
    pipeline;
    \item verifies that target and teacher have identical planner interface
    values and identical canonical environment JSON;
    \item executes the teacher controller for its configured learning and
    evaluation cycles;
    \item records realized transitions after each action;
    \item rejects the run if fewer than \(H\) transitions were produced;
    \item builds the target standalone teacher contract;
    \item validates the resulting artifact by constructing the target
    \Warm{} agent from it;
    \item writes pretty canonical teacher JSON to \texttt{--out}.
\end{enumerate}
This realizes the offline-to-online path: any controller can act as teacher,
and the resulting artifact initializes a \Warm{} student for the target task.

The reproducibility CLI also exposes:
\begin{verbatim}
infotheory warmstart teacher from-jsonl \
  --target <warmstart-planner-run> --jsonl <run.jsonl> --out <teacher.json>

infotheory warmstart teacher merge \
  --target <warmstart-planner-run> --out <teacher.json> \
  --teacher <teacher-a.json> --teacher <teacher-b.json> ...
\end{verbatim}
The first command implements the deterministic telemetry converter from
Section~\ref{sec:telemetry}. The second performs the deterministic
same-contract trace union defined there and validates the merged result by
constructing the target student.

\section{Online Trace Refresh}
\label{sec:refresh}

During online deployment, \Warm{} can export its own live same-task trace after
at least \(H\) transitions. The live trace is not itself a teacher artifact; it
becomes teacher data only when it is wrapped under the unchanged target
contract, merged into the teacher dataset, and validated by the same full
dataset gate used for ordinary teacher artifacts. A refresh wrapper may then
reinitialize the student from the merged dataset.

A bounded refresh loop is specified by an initial teacher dataset
\(\Teach^{(0)}\), a maximum number of refresh rounds \(R\), and a finite run
budget per round. Round \(r\) builds \(W^{(r)}\) from \(\Teach^{(r)}\), executes
it on the same compiled bridge, extracts a live trace \(\tau^{(r)}\) if one has
at least \(H\) transitions, and sets
\[
\Teach^{(r+1)}
=
\operatorname{Merge}(\Teach^{(r)},\tau^{(r)}).
\]
\(\operatorname{Merge}\) is the deterministic trace-content union of
Section~\ref{sec:teacher-artifact}. The loop stops after \(R\) rounds, or
earlier if no admissible live trace is produced.

\begin{proposition}[Refresh trace payload admissibility]
If a live \Warm{} trace of length at least \(H\) is produced by a target
runtime that validates every emitted action, observation, and reward against
the compiled target bridge, then wrapping that trace under the unchanged target
contract yields transition payloads whose complete horizon windows are
admissible teacher labels for that same target contract.
\end{proposition}

\begin{proof}
Every recorded action, observation, and reward has already passed the target
runtime's transition validator. Each label is the deterministic \(H\)-step sum
of rewards in that same trace. The trace length premise gives at least one
complete delayed-label window, and the unchanged wrapper contract supplies the
same fingerprint, interface fields, horizon, phase period, and provenance
fields required by Definition~\ref{def:admissible-teacher}. Therefore the
wrapped live trace can pass the teacher-artifact validator as part of the
merged dataset.
\end{proof}

Refresh is a data mechanism, not a monotonic-improvement theorem. It improves
the theorem-facing local bound only to the extent that the refreshed predictor
is closer to the target conditional label law on visited histories.

\begin{theorem}[Conditional local refresh improvement]
Fix a history \(h\) and suppose two same-bridge students differ only in their
teacher artifact: an initial student \(W_0\) and a refreshed student \(W_1\).
If, for every action,
\[
\|\widehat p_{W_1}(\cdot\mid h,a)-p(\cdot\mid h,a)\|_{\TV}
<
\|\widehat p_{W_0}(\cdot\mid h,a)-p(\cdot\mid h,a)\|_{\TV},
\]
then the local finite-horizon value-error bound of
Theorem~\ref{thm:tv-robustness} is strictly smaller for \(W_1\) than for
\(W_0\) at \(h\).
\end{theorem}

\begin{proof}
The bound in Theorem~\ref{thm:tv-robustness} is monotone in the uniform
total-variation radius. A strict decrease for every action gives a strict
decrease of the supremum radius at \(h\), and therefore a strict decrease of
the displayed value-error bound.
\end{proof}

For bridges with monotone incumbent semantics, such as the tuner exact-\(\Jh\)
bridge, the bridge may prove additional best-so-far or objective-monotonicity
claims. Those claims belong to the bridge specification; they are not implied
by refresh alone for arbitrary environments.

\section{Core Theorems}
\label{sec:theorems}

Let \(p_t(\ell\mid h,a)\) be the true conditional law of the exact
\(H\)-step return label at history \(h\) under action \(a\). Define
\[
Q^\star_H(h,a)=\sum_{\ell\in\Lab}\Vlab(\ell)p_t(\ell\mid h,a).
\]
Let
\[
D_G=\max_{\ell\in\Lab}\Vlab(\ell)-\min_{\ell\in\Lab}\Vlab(\ell).
\]
For the direct integer-reward bridge, \(D_G=Hr_{\max}-Hr_{\min}=Hr_{\max}\).

\begin{theorem}[Exact greedy alignment]
If, at a greedy decision history \(h\),
\[
\widehat p_t(\cdot\mid h,a)=p_t(\cdot\mid h,a)
\quad\text{for every }a\in\Act,
\]
then \Warm{} selects an action maximizing the true expected exact
\(H\)-step return \(Q^\star_H(h,a)\), with deterministic lowest-index
tie-breaking.
\end{theorem}

\begin{proof}
Under the premise, \(\widehat Q_H(h,a)=Q^\star_H(h,a)\) for every action.
The implemented rule chooses the lowest-index maximizer of these same values.
\end{proof}

\begin{theorem}[Total-variation robustness]
\label{thm:tv-robustness}
If
\[
\sup_{a\in\Act}
\|\widehat p_t(\cdot\mid h,a)-p_t(\cdot\mid h,a)\|_{\TV}
\le \delta,
\]
then the greedy \Warm{} action \(\hat a\) satisfies
\[
Q^\star_H(h,a^\star)-Q^\star_H(h,\hat a)\le 2D_G\delta,
\]
where \(a^\star\) is any true optimal action.
\end{theorem}

\begin{proof}
For any bounded function with range diameter \(D_G\), expectation changes by
at most \(D_G\) times total variation. Thus
\[
|\widehat Q_H(h,a)-Q^\star_H(h,a)|\le D_G\delta
\]
for every action. Since \(\hat a\) maximizes \(\widehat Q_H\),
\[
Q^\star_H(h,a^\star)
\le \widehat Q_H(h,a^\star)+D_G\delta
\le \widehat Q_H(h,\hat a)+D_G\delta
\le Q^\star_H(h,\hat a)+2D_G\delta.
\]
\end{proof}

\begin{corollary}[Teacher action-gap preservation]
If the teacher value \(\QT\) has a unique greedy action \(a_T\) at \(h\) and
\[
\sup_{a\in\Act}|\widehat Q_W(h,a)-\QT(h,a)|<\Delta_T(h)/2,
\]
then the greedy warmstart action equals \(a_T\). In particular, the condition
holds whenever
\[
\sup_{a\in\Act}\|\widehat p_W(\cdot\mid h,a)-\PT(\cdot\mid h,a)\|_{\TV}
\le\delta
\quad\text{and}\quad
D_G\delta<\Delta_T(h)/2.
\]
\end{corollary}

\section{Tuner Bridge Instantiation}
\label{sec:tuner-bridge}

The tuner bridge instantiates the generic reward \(r_t\) as exact improvement
in the deployable two-part objective:
\[
\Jh(z;d)=8L_B(z)+\ell_H(z).
\]
For incumbent sequence \(b_t\), the exact reward is
\[
r_t=\Jh(b_{t-1};d)-\Jh(b_t;d)
\]
after the bridge's finite reward encoder has certified representability and
injectivity. Therefore the \(H\)-step label telescopes:
\[
G_i^{(H)}
=\sum_{j=i}^{i+H-1}r_j
=\Jh(b_{i-1};d)-\Jh(b_{i+H-1};d).
\]
This is the tuner theorem use case. Its dataset lowering, deployability,
timing, exact reward certificate, and observation adapter are governed by the
tuner specification. \Warm{} consumes the resulting externally validated
teacher artifact through the same contract fields, but this standalone
document does not weaken or replace the tuner theorem assumptions.

\section{Implementation Conformance}
\label{sec:implementation}

A conforming implementation must:
\begin{itemize}
    \item route planner-run documents through the canonical spec
    validate/compile pipeline;
    \item reject zero return horizon, zero return bins, phase period below
    horizon, planner simulation budgets other than the direct-evaluator marker
    \(1\), unsupported predictor conditioning, and malformed teacher contracts;
    \item validate every teacher and live transition against action,
    observation, and reward bounds before learning from it;
    \item use deterministic LSB-first fixed-width field encodings for actions,
    observations, and rewards, and MSB-first value-monotone label codewords;
    \item reset transient predictor conditioning history between independent
    teacher traces;
    \item expose canonical teacher JSON serialization and parsing as inverse
    operations on valid artifacts;
    \item keep JSONL telemetry separate from canonical teacher artifacts;
    \item merge same-contract traces by deterministic content deduplication;
    \item record whether actions were greedy or exploratory in planner JSONL
    telemetry when provenance is present.
\end{itemize}

\section{Traceability Obligations}
\label{sec:traceability}

The implementation and tests should maintain the following traceability map:
\begin{itemize}
    \item compiled bridge contract: parser, canonical JSON, binary round-trip,
    compile, and fingerprint tests for \texttt{aiqi\_warmstart\_exact\_jh};
    \item admissible teacher artifacts: schema parse/serialize inverse tests,
    contract mismatch rejection tests, standalone planner-run export tests, and
    tuner-bridge validation tests owned by the tuner specification;
    \item JSONL conversion: same-step action/percept pairing tests,
    pre-action percept next-step pairing tests, mixed-convention rejection
    tests, malformed/duplicate rejection tests, and final target-agent
    construction from converted data;
    \item deterministic merge and refresh: structural trace-content deduplication
    tests, same-contract validation tests, and rebuilt-agent construction from
    merged data;
    \item online delayed labels: within-run transition observation tests showing
    a newly closed \(H\)-window updates exactly one phase model;
    \item implementation anchors:
\begin{verbatim}
PlannerActionProvenance
warmstart_teacher_trace_from_jsonl_reader
warmstart_teacher_trace_from_jsonl_path
merge_warmstart_teacher_traces_deterministic
standalone_warmstart_teacher_contract_for_compiled_planner_run
warmstart_exact_jh_planner_task_fingerprint
WarmStartExactJhAgent::from_compiled_planner_run
validate_warmstart_exact_jh_teacher_contract
load_warmstart_exact_jh_teacher_dataset
PlannerControllerAgent::from_compiled
PlannerEnvironment
PlannerRunSession
run_warmstart_teacher_planner_run_export
run_warmstart_teacher_from_jsonl_export
run_warmstart_teacher_merge
run_compiled_planner_run
\end{verbatim}
    \item theorem claims: total-variation/action-gap claims are mathematical
    obligations; implementation tests must verify their preconditions are
    exposed and not silently inferred from unvalidated data.
\end{itemize}

\section{Claim Boundary}
\label{sec:claims}

\Warm{} does not claim universal optimality, discounted-return optimality,
wall-clock dominance over its teacher, or automatic repair of missing
counterfactual action coverage. The exact theorem is local and finite-horizon:
if the conditional label law is exact at the current history, greedy
\Warm{} is exact for the current \(H\)-step return objective. Offline teacher
data, online refresh, and exploration are mechanisms for making that local
model useful and reproducible; their stronger empirical claims must be
measured and reported as experimental results.

\end{document}
